\documentclass[11pt, reqno]{amsart}

\usepackage[spanish]{babel}
\usepackage[LGR, T1]{fontenc}
\usepackage[utf8]{inputenc}

../LaTeX/general.tex
% \input{../graphics.tex}

\makeatletter
\def\emailaddrname{\textit{Correo electrónico}}
\def\subtitle#1{\gdef\@subtitle{#1}}
\def\@subtitle{}

% Metadata
\def\logo#1{\gdef\@logo{#1}}
\def\@logo{}
\def\institution#1{\gdef\@institution{#1}}
\def\@institution{}
\def\department#1{\gdef\@department{#1}}
\def\@department{}
\def\professor#1{\gdef\@professor{#1}}
\def\@professor{}
\def\course#1{\gdef\@course{#1}}
\def\@course{}
\def\coursecode#1{\gdef\@coursecode{#1}}
\def\@coursecode{}

\renewcommand{\maketitle}{
\begin{center}
	\small
	\renewcommand{\arraystretch}{1.2}
	\begin{tabular}{cp{.37\textwidth}p{0.44\textwidth}}
		% \hline
		\multirow{5}{*}{\includegraphics[height=2.0cm]{\@logo}}
	  & \multicolumn{2}{c}{ \makecell{{\bfseries \@institution} \\ \@department} } \\
	  % & \multicolumn{2}{c|}{{\bfseries\@institution} \\ \@department} \\
	  \cline{2-3}
	  & \textbf{Profesor:} \@professor & \textbf{Ayudante:} \authors \\
	  % \cline{2-3}
	  & \textbf{Curso:} \@course & \textbf{Sigla:} \@coursecode \\
	  % \cline{2-3}
	  & \multicolumn{2}{l}{ \textbf{Fecha:} \@date } \\
	  % \hline
	\end{tabular}
	\\[\baselineskip]
	% {}
	% \vspace{2\baselineskip}
	{\bfseries\Large\@title}
	\ifx\@subtitle\@empty\else
		\\[1ex]
		\large\mdseries\@subtitle
	\fi
\end{center}
}
\makeatother

\usepackage{multirow, makecell}

\usepackage[
	reversemp,
	letterpaper,
	% marginpar=2cm,
	% marginsep=1pt,
	margin=2.3cm
]{geometry}
\usepackage{fontawesome}
% \makeatletter
% \@reversemargintrue
% \makeatother

% Símbolos al margen, necesitan doble compilación
\newcommand{\hard}{\marginnote{\faFire}}
\newcommand{\hhard}{\marginnote{\faFire\faFire}}

% Dependencias para los teoremas
\usepackage{xifthen}
\def\@thmdep{}
\newcommand{\thmdep}[1]{
	\ifthenelse{\isempty{#1}}
	{\def\@thmdep{}}
	{\def\@thmdep{ (#1)}}
}
\newcommand{\thmstyle}{\color{thm}\sffamily\bfseries}

% ===== Estilos de Teoremas ==========
\newtheoremstyle{axiomstyle}
	{0.3cm}
	{0.3cm}
	{\normalfont}
	{0.5cm}
	{\bfseries\scshape}
	{:}
	{4pt}
	{\thmname{#1}\thmnote{ #3}\thmnumber{ (#2)}}
\newtheoremstyle{styleC}
	{0.5cm}
	{0.5cm}
	{\normalfont}
	{0.5cm}
	{\bfseries}
	{:}
	{4pt}
	{\thmname{#1\textrm{\@thmdep}}\thmnumber{ #2}\thmnote{ (#3)}}

% ====== Teoremas (sin borde) ===========
\theoremstyle{axiomstyle}
\newtheorem*{axiom}{Axioma}

% ====== Teoremas (sin borde) ==================
\theoremstyle{styleC}
\newtheorem{thm}{Teorema}[section]
\newtheorem{mydef}[thm]{Definición}
\newtheorem{prop}[thm]{Proposición}
\newtheorem{cor}[thm]{Corolario}
\newtheorem{lem}[thm]{Lema}
\newtheorem{con}[thm]{Conjetura}
\newtheorem*{sol}{Solución}
\newtheorem*{obs}{Observación}
\newtheorem*{ex}{Ejemplo}

\newtcbox{bluebox}[1][]{enhanced jigsaw, 
  sharp corners,
  frame hidden,
  nobeforeafter,
  listing only,
  #1} % comando para crear cajas de colores

\expandafter\let\expandafter\oldproof\csname\string\proof\endcsname
\let\oldendproof\endproof
\renewenvironment{proof}[1][\proofname]{%
  \oldproof[\scshape Demostración:]%
}{\oldendproof} % comando para redefinir la caja de la demostración
\newenvironment{hint}[1][\proofname]{%
  \oldproof[\scshape Pista:]%
}{\oldendproof} % comando para redefinir la caja de la demostración

% colores utilizados
\definecolor{numchap}{RGB}{249,133,29}
\definecolor{chap}{RGB}{6,129,204}
\definecolor{sec}{RGB}{204,0,0}
\definecolor{thm}{RGB}{106,176,240}
\definecolor{thmB}{RGB}{32,31,31}
\definecolor{part}{RGB}{212,66,66}

% ====== Diseño de los titulares ===============
\usepackage[explicit]{titlesec} % para personalizar el documento, la opción <<explicit>> hace que el texto de los titulares sea un objeto interactuable

\titleformat{\subsection}[runin]
	{\bfseries}
	{\textrm{\S}\thesubsection}
	{1ex}
	{#1.}

\setlist[enumerate,1]{label=\arabic*., ref=\arabic*} % Enumerate standards

\usepackage{tikz}
\usetikzlibrary{babel,cd}
% \DeclareMathOperator\Gr{Gr}

% \DeclareMathOperator\Cono{Cono}

\title{El teorema del punto fijo de Hopf--Lefschetz}
\date{\DTMdate{2026-09-24}}

\author{José Cuevas Barrientos}
\email{josecuevasbtos@uc.cl}
\urladdr{https://josecuevas.xyz/teach/2026-2-ayud/}

\logo{../puc_negro.png}
\institution{Pontificia Universidad Católica de Chile}
\department{Facultad de Matemáticas}
\course{Topología Algebraica}
\coursecode{MAT2850}
\professor{Mauricio Bustamante}

\begin{document}

\maketitle

Recuerde que si $V$ es un espacio vectorial de dimensión finita sobre un cuerpo arbitrario y $\varphi \in \End(V)$ es un endomorfismo,
entonces la \emph{traza} $\tr\varphi$ es la traza usual de cualquier representación matricial de $\varphi$ (i.e., la suma de su diagonal).
No depende de la elección de base ordenada por la fórmula usual $\tr(A\cdot B) = \tr(B\cdot A)$.

\section{Traza}
\begin{enumerate}
	\item Sea $f \colon V \to V$ un endomorfismo de un espacio vectorial y sea $W \le V$ un subespacio $\varphi$-estable (i.e., $\varphi[W] \subseteq W$).
		Pruebe que $\tr\varphi = \tr(\varphi|_W) + \tr(\overline{\varphi} \colon V/W \to V/W)$.

	\item \textbf{Fórmula de la traza de Hopf:}
		Sea $C_\bullet$ un complejo finito de espacios vectoriales de dimensión finita y sea $\varphi_\bullet \colon C_\bullet \to
		C_\bullet$ un endomorfismo.
		Pruebe que
		\[
			\sum_{n\in\Z} (-1)^n \tr(\varphi_n) = \sum_{n\in\Z} (-1)^n \tr H_n(\varphi).
		\]
\end{enumerate}

\section{El teorema de Lefschetz}
En esta sección, $H_n(K)_\Q := H_n(K; \Z) \otimes_\Z \Q$.
Cuando $K$ es un complejo finito, recuerde que los grupos de homología $H_n(K; \Q)$ son finitamente generados, de modo que los $H_n(K)_\Q$
son $\Q$-espacios vectoriales de dimensión finita.

\begin{enumerate}[resume]
	\item Sea $K$ un complejo simplicial finito y $f \colon |K| \to |K|$ una función continua.
		Defina
		\[
			\mathcal{L}(f) := \sum_{n\ge 0} (-1)^n \tr H_n(f)_\Q.
		\]
		\begin{enumerate}
			\item Pruebe que $\mathcal{L}(\Id) = \chi(K)$.

			\item Pruebe que $\mathcal{L}(f)$ es invariante bajo homotopía.

			\item \textbf{Teorema del punto fijo de Lefschetz:}
				Demuestre que si $f$ no tiene puntos fijos, entonces $\mathcal{L}(f) = 0$.

				\begin{hint}
					Para esto, considere una aproximación simplicial a $f$ que, como bien sabe el lector, requiere
					subdividir primero al complejo $K$.
					Tras encajar a la realización geométrica $|K| \subseteq \R^N$ en un espacio euclídeo, pruebe que
					\[
						\delta := \min\{ |f(x) - x| : x \in |K| \} > 0.
					\]
					Mediante subdivisión baricéntrica puede suponer que la subdivisión es tal que las caras tienen
					diámetro menor que $\delta$, ¿qué sucede entonces con la aproximación simplicial?
				\end{hint}
		\end{enumerate}

	\item Concluya el \textbf{teorema del punto fijo de Brouwer:}
		toda función continua $f$ sobre la $n$-bola cerrada $\D^n$ tiene un punto fijo.
	\item Calcule $\chi(\Sn)$.
\end{enumerate}

\appendix
\section{Comentarios}
El lector instruído en la teoría de variedades diferenciales podría intentar demostrar las siguientes consecuencias al teorema de Lefschetz
(o en su defecto consultar \citeauthor{fomenko:homotopical}~\cite[235]{fomenko:homotopical}):
\begin{enumerate}[resume]
	\item Sea $X$ una variedad diferencial con borde compacta, y sea $\xi$ un campo vectorial sobre $X$.
		Suponga que $\xi$ no se anula en ningún punto y que está dirigida en el borde $\partial X$.
		Pruebe que $\chi(X) = 0$.

	\item \textbf{El \textquote{teorema de la bola peluda}:}
		Pruebe que todo campo vectorial sobre $\Sn[2]$ se anula en al menos un punto.
\end{enumerate}

El número de Lefschetz $\mathcal{L}(f)$ no funciona \textquote{porque sí}.
Tenemos la siguiente caracterización algebraica:
\begin{thm}
	Sea $X$ una $n$-variedad compleja compacta (e.g., proyectiva) y sea $f \colon X \to X$ un morfismo de $\C$-variedades.
	Si $\Gamma_f, \Delta$ denotan las subvariedades de $X \times X$ de dimensión $n$ correspondientes al gráfico de $f$ y a la diagonal
	$\Delta$ resp., entonces el número de intersección
	\[
		(\Gamma_f.\Delta) = \mathcal{L}(f) = \sum_{q\ge 0} (-1)^q \tr H^q(f)_\Q,
	\]
	donde $H^q(X)_\Q = H_{\rm sing}^q(X; \Q)$ denota el grupo de cohomología singular (más usual en geometría algebraica).
\end{thm}
En particular, si la intersección de la diagonal es transversal y \textquote{orientada} (en el sentido de que $\det(I_n - \ud_x f) > 0$ para
todo punto fijo $x$), entonces $\mathcal{L}(f)$ cuenta exactamente la cantidad de puntos fijos.
% Otra consecuencia es que $f$ tiene al menos $|\mathcal{L}(f)|$

Definiendo \textquote{número de intersección} (cfr.\ \cite{fomenko:homotopical}, \S 17.5) para variedades diferenciales cerradas (i.e.,
compactas sin borde) orientables y triángulables (i.e., homeomorfas a un $|K|$ para un complejo finito $K$), el teorema anterior sigue
siendo válido.

\printbibliography

\end{document}
